
\chapter{Literature Review}\label{chap:litreview}

\section{Review Scope and Questions}
This chapter surveys work relevant to securing voice assistants against spoofing, with emphasis on: (i) attacks and defenses in voice interfaces; (ii) audio watermarking for speech provenance; and (iii) deepfake/replay detection. The literature review intentionally covers both the \emph{implemented} watermarking branch of the current thesis and the \emph{planned} user-side detector branch, because the latter remains important for the long-term system architecture even though it is not experimentally realized in the present prototype.

\section{Review Method (Concise Summary)}
We adopt a structured narrative review aligned with established software-engineering evidence guidelines. At a high level, we (1) defined broad questions and a taxonomy; (2) searched major digital libraries and reputable preprint servers using terms such as \emph{voice assistant security}, \emph{audio watermark*}, \emph{speech deepfake}, \emph{spoofing}, \emph{replay detection}, and \emph{ASVspoof}; (3) prioritized peer-reviewed papers and widely cited technical reports; and (4) synthesized findings by theme. Because our goal is to scope and synthesize (not to conduct a full SLR/meta-analysis), we report a concise method summary and do not claim exhaustiveness; however, we follow the spirit of systematic mapping---transparent scope, inclusion logic, and categorized synthesis---commonly used in CS reviews.

\section{Taxonomy Used in This Review}
We organize the field into four themes that map to the dual-channel (user$\rightarrow$agent / agent$\rightarrow$user) nature of conversational systems:
\begin{enumerate}[leftmargin=1.25em]
  \item \textbf{Threats and Attacks in Voice Interfaces} (replay, TTS/VC deepfakes, injection vectors).
  \item \textbf{Audio Watermarking for Speech Provenance} (payload and zero-bit watermarks; robustness vs.\ imperceptibility; deployment in TTS).
  \item \textbf{Deepfake and Replay Detection} (features, end-to-end models, generalization and calibration).
  \item \textbf{Hybrid/Integrated Defenses} (joint watermarking--detection and attribution).
\end{enumerate}

\section{Voice Assistant Security and AI-Generated Audio Threats}
Voice-based digital assistants pose unique security and privacy challenges due to their open-microphone design and reliance on audio commands\,[15]. Prior work documents replay attacks (recorded voice commands or wake-words) and, more recently, high-quality AI voice cloning. With modern TTS/VC, a few seconds of someone’s voice---or even cross-modal cues such as photos in specialized systems like Foice---can suffice to synthesize arbitrary speech in their timbre\,[3][16]. This enables impersonation and vishing at scale. Public evaluations have also shown that commercial voice authentication can be fooled by cloned voices\,[3]. The ASV literature further shows that, without countermeasures, automatic speaker verification is vulnerable to both replay and synthetic speech\,[17][18].

\paragraph{Attack directions.}
\textbf{(U$\rightarrow$A)} An adversary impersonates a user or injects commands via replay or TTS/VC, relying on ASR and any voice verification to accept the content. \textbf{(A$\rightarrow$U)} A fake voice pretends to be the assistant (e.g., in phishing calls or man-in-the-middle tampering). While the first direction is well studied, the second motivates watermarking for assistant speech so verifiers can validate provenance.

\section{Audio Watermarking for Speech Provenance}
Audio watermarking embeds a secret pattern that is inaudible to humans but detectable by machine, even after typical modifications. Within AI speech, this provides provenance for synthetic audio\,[21][22]. We focus on \emph{proactive} watermarking: the assistant marks its own TTS output so that receivers can verify authenticity.

\paragraph{Representative approaches.}
Zong et al.\ (AudioMarkNet, 2025) embed watermarks into publicly released recordings so that a future cloned voice unwittingly reproduces the mark, enabling simple verification\,[23][24]. Fernandez et al.\ (AudioSeal, 2024) jointly train an embedder and a fast, localized detector with perceptual masking to keep changes inaudible and resilient to cropping/encoding/noise\,[25][26][27][14]. Other works balance fidelity and robustness via redundancy or augmentation (training with distortions and codecs)\,[22][28]. Liu et al.\ (2024) propose timbre watermarking using an encoder--decoder; while highly hidden, it can be brittle under noise+resampling\,[29][30][31]. In contrast, Zong et al.\ report simpler, repeated patterns that survived such attacks when timbre watermarking failed\,[12]. Li et al.\ (TRUE, 2025) emphasize time-domain learning to preserve fidelity while remaining robust, achieving near-original PESQ\,[32][33][34]. For our use case, zero-bit watermarking (presence/absence) is often attractive\,[35][36][37], as it minimizes payload and favors robustness; however, the current implemented prototype uses a \textbf{small 4-bit payload} rather than a strictly zero-bit mark.

\section{Deepfake and Replay Audio Detection}
Audio deepfake detection distinguishes genuine human speech from synthesized or tampered audio (including replay). Early systems used hand-crafted features (MFCC, CQCC) with classical classifiers; recent work favors deep models on spectrograms or raw waveforms, sometimes leveraging large pre-trained speech encoders\,[38][39][40][5]. A persistent challenge is \emph{generalization}: detectors often excel on known attacks but degrade on unseen methods or channel conditions\,[41][42][6][8]. Reports show EERs rising above 20\% on out-of-distribution data\,[43]. Robustness strategies include diverse attack/channel augmentation, model ensembles, and calibration using DET curves/EER to set operating thresholds\,[44][45][13]. Replay detection additionally exploits environmental cues (microphone/speaker colorations, reverberation), where deep models can learn discriminative patterns when trained with realistic replay data.

\paragraph{Impact of the transmission channel.}
Assistants frequently receive compressed and noisy audio; codecs and room acoustics can blur artifacts. Some findings indicate artifact-based detection may remain informative under neural codecs while naïve watermarking may be stripped\,[50]. This trade-off motivates combining detection with watermarking.

\section{Hybrid/Integrated Defenses}
Hybrid designs strengthen resilience by offering multiple, complementary signals. FakeMark (2025) injects watermarked artifacts that aid attribution and remain useful even if one cue is weakened\,[51][52][53][54]. This dual-cue principle underpins the \emph{long-term architecture} of our system vision: attackers would ideally need to both (i) evade or remove the assistant’s watermark and (ii) leave no detectable synthesis/replay traces. However, in the current thesis this second detector cue remains conceptual rather than implemented, so the empirical work focuses on the watermarking branch alone.

\section{Synthesis: What the Literature Implies for Our Design}
\textbf{(S1) Provenance works in practice.} Modern speech watermarks can be imperceptible (MOS-comparable to unwatermarked TTS) yet resilient to real-world edits\,[14][25][27]. \textbf{(S2) Detection remains necessary.} Generalization is hard; unseen generators and channels inflate EER\,[6][8][41][43]. \textbf{(S3) Layering helps.} Combining watermarking with artifact detection hardens systems against both removal attempts and generator evolution\,[51][52]. \textbf{(S4) Channel realism is critical.} Training/evaluating under realistic distortions and codecs improves robustness and yields honest operating points\,[44][45][13].

\subsection{Our design vs.\ prior art}
The current thesis differs from prior work in several concrete ways:
\begin{itemize}[leftmargin=1.25em]
  \item \textbf{Implemented scope is narrower than the full architecture.} Unlike hybrid proposals such as FakeMark that are discussed as integrated multi-cue defenses, the present prototype \emph{implements only the assistant-output watermarking branch}; the user-side deepfake/replay detector is an integration path for future work.
  \item \textbf{Localized, prefix-aware verification is a primary design target.} Relative to AudioSeal-style localized detection, our system places unusual emphasis on whether the watermark can already be \emph{decoded from short prefixes} (0.25/0.5/1.0/2.0~s), not merely detected after the full utterance.
  \item \textbf{Small-payload watermark instead of strict zero-bit detection.} Whereas many provenance systems emphasize presence/absence only, the current prototype embeds a \textbf{4-bit payload} and decodes it through masked statistics pooling over the predicted watermark region.
  \item \textbf{Training objective prioritizes functionality before perceptual refinement.} In contrast to systems such as TRUE that strongly emphasize fidelity, the stable checkpoints in this thesis first prioritize reliable detection, localization, and early decoding; quality-oriented tuning is explicitly deferred to the next engineering stage.
  \item \textbf{Evaluation scope is intentionally constrained.} Rather than claiming broad cross-domain or full VoIP-room-acoustics generalization, the implemented benchmark is limited to \textbf{LibriSpeech} and a concrete black-box attack suite (noise, speed/resampling, low-pass, crop-pad, bit-depth, MP3, AAC), which makes the claims narrower but better aligned with the actual evidence.
\end{itemize}

\section{Gaps and Research Opportunities}
\begin{itemize}[leftmargin=1.25em]
  \item \textbf{Watermark removal under targeted attack.} Systematic studies of adversarial removal vs.\ imperceptibility are still limited, especially for black-box cloners.
  \item \textbf{Cross-model generalization.} Detectors need stronger transfer to unseen TTS/VC families and neural codecs without frequent retraining.
  \item \textbf{End-to-end V2V evaluations.} Few works report both provenance \emph{and} detection outcomes with interactive latency and user-perceived quality.
\end{itemize}
These gaps motivate the present thesis in two ways: first, by focusing the implemented work on a reproducible watermarking benchmark; and second, by keeping the user-side detector and full bilateral middleware as clearly labeled future extensions rather than overstating what has already been built.

\section{Threats to the Validity of This Review}
\textbf{Construct validity.} Our taxonomy focuses on watermarking and detection for voice assistants; broader audio forensics is only tangentially covered. \textbf{Internal validity.} Because this is a structured narrative review rather than a full SLR, selection bias is possible; we mitigated it by searching multiple venues and preferring peer-reviewed and widely cited sources. \textbf{External validity.} Findings may not generalize to all languages, devices, or proprietary codecs; we therefore prioritize design patterns (layering, channel realism) over narrow system choices.
